% Part: normal-modal-logic
% Chapter: frame-definability
% Section: introduction

\documentclass[../../../include/open-logic-section]{subfiles}

\begin{document}

\olfileid{nml}{frd}{int}

\olsection{પરિચય}

મોડલ તર્કશાસ્ત્રીઓને રસ પડે એવો એક પ્રશ્ન સુલભતા
સંબંધ અને તે સંબંધ ધરાવતા નિદર્શનોમાં ચોક્કસ
!!{formula}sની સત્યતા વચ્ચેના સંબંધનો છે. ઉદાહરણ
તરીકે, માનો કે સુલભતા સંબંધ સ્વવાચક છે, એટલે કે
દરેક $w
\in W$ માટે $Rww$. બીજા શબ્દોમાં, દરેક જગત
પોતામાંથી સુલભ છે. તેનો અર્થ એ કે કોઈ જગત~$w$માં
$\Box !A$ સત્ય હોય ત્યારે $w$ પોતે પણ એવા સુલભ
જગતોમાંનું એક છે જ્યાં~$!A$ સત્ય હોવું પડે. તેથી
$\mModel{M}$નો સુલભતા સંબંધ~$R$ સ્વવાચક હોય, તો
કોઈપણ જગત $w$ અને સૂત્ર $!A$ લઈએ, ત્યાં $\Box !A
\lif !A$ સત્ય છે (બીજા શબ્દોમાં, સૂત્ર-ઢાંચો $\Box p \lif
p$ અને તેના બધા પ્રતિસ્થાપન દૃષ્ટાંતો~$\mModel{M}$માં
સત્ય છે).

પરંતુ વિપરીત વિધાન ખોટું છે. દાખલા તરીકે,
$\Box p \lif p$ એ~$\mModel{M}$માં સત્ય હોય તો $R$
સ્વવાચક છે એવું નથી. એક જ જગત~$w$ ધરાવતું એવું
અસ્વવાચક નિદર્શન~$\mModel{M}$ સરળતાથી મળે જેમાં
$\Box p \lif
p$ દરેક જગતમાં સત્ય છે: તે પોતામાંથી સુલભ
ન હોય, પણ $w \in V(p)$ હોય એવું નિદર્શન લો.
$p$નું સત્યમૂલ્ય યોગ્ય રીતે પસંદ કરીને અસ્વવાચક
નિદર્શનમાં $\Box !A \lif !A$ સત્ય બનાવી શકીએ છીએ.

ઉપાય એ છે કે સમીકરણમાંથી ચલ-નિમણૂક~$V$ને દૂર કરીએ.
જો~$\mModel{M}$ના બધા જગતોમાં $\Box p \lif p$ સત્ય
હોવાની માગણી~$V(p)$માં કયાં જગતો છે તેનાથી સ્વતંત્ર
રીતે કરીએ, તો $R$ સ્વવાચક હોવું જરૂરી છે. કેમ કે
કોઈપણ અસ્વવાચક નિદર્શનમાં ઓછામાં ઓછું એક જગત~$w$
એવું હશે કે $Rww$ ન હોય. જો $V(p) =
W \setminus \{w\}$ મૂકીએ, તો $p$ એ~$w$ સિવાયનાં
બધાં જગતોમાં, તેથી~$w$માંથી સુલભ બધા જગતોમાં,
સત્ય હશે ($w$ એ~$w$માંથી સુલભ નથી અને $p$ માત્ર
$w$માં જ અસત્ય છે). બીજી બાજુ, $p$ એ $w$માં
અસત્ય છે, તેથી $\Box
p \lif p$ એ~$w$માં અસત્ય છે.

આથી નિમણૂક વિનાની નિદર્શન-સંરચનાઓ માટે સંકેત
રજૂ કરવો યોગ્ય લાગે છે: તેમને \emph{ફ્રેમ} કહીએ.
ફ્રેમ $\mModel{F}$ એ માત્ર જગતોના ગણ અને તેમના
સુલભતા સંબંધની જોડ $\tuple{W, R}$ છે. ત્યારે
દરેક નિદર્શન $\tuple{W, R, V}$ આ ફ્રેમ
$\tuple{W, R}$ \emph{પર આધારિત} છે એમ કહીએ.
વિપરીત રીતે, ફ્રેમ તેના પર આધારિત નિદર્શનોનો
વર્ગ નક્કી કરે છે; અને ફ્રેમોના વર્ગથી તે વર્ગની
કોઈપણ ફ્રેમ પર આધારિત નિદર્શનોનો વર્ગ નક્કી થાય છે.
વળી, ફ્રેમમાં !!a{formula}ની \emph{માન્યતા} માટે
$\mModel{F} \Entails !A$ની વ્યાખ્યા આપી શકીએ: આનો
અર્થ~$\mModel{F}$ પર આધારિત દરેક $\mModel{M}$ માટે
$\mSat{M}{!A}$ છે.

આ સંકેતથી !!{formula}s અને ફ્રેમોના વર્ગો વચ્ચે
અનુરૂપતા સંબંધો સ્થાપી શકીએ: દાખલા તરીકે,
$\mModel{F} \Entails \Box p \lif p$ ત્યારે અને માત્ર
ત્યારે જ જ્યારે $\mModel{F}$ સ્વવાચક હોય.

\end{document}
